% Figura: problema 5.2. Dipolo puntual en el centro de una esfera de radio a:
% líneas de campo exteriores (exactas) y distribución angular de la densidad de energía.
\begin{figure}[t]
\centering
\begin{subfigure}[b]{0.52\textwidth}\centering
\begin{tikzpicture}
\begin{axis}[marco, width=7.2cm, xmin=-2.6, xmax=2.6, ymin=-2.0, ymax=2.0]
  \fill[dielec] (axis cs:0,0) circle[radius=1];
  %% Generado por scripts/tema5.py -- no editar a mano
\addplot[lineaE, flechaen=0.396] coordinates {(2.600,1.454) (2.571,1.426) (2.542,1.398) (2.513,1.370) (2.484,1.342) (2.455,1.314) (2.425,1.286) (2.395,1.259) (2.365,1.232) (2.335,1.205) (2.305,1.178) (2.274,1.151) (2.244,1.125) (2.213,1.098) (2.182,1.072) (2.151,1.047) (2.120,1.021) (2.089,0.995) (2.057,0.970) (2.025,0.945) (1.993,0.920) (1.961,0.895) (1.929,0.871) (1.897,0.847) (1.864,0.823) (1.831,0.799) (1.799,0.775) (1.766,0.752) (1.732,0.729) (1.699,0.706) (1.666,0.683) (1.632,0.661) (1.598,0.639) (1.564,0.617) (1.530,0.595) (1.496,0.573) (1.461,0.552) (1.427,0.531) (1.392,0.511) (1.357,0.490) (1.322,0.470) (1.287,0.450) (1.251,0.431) (1.216,0.411) (1.180,0.392) (1.144,0.374) (1.108,0.355) (1.072,0.337) (1.036,0.319) (0.999,0.302) (0.963,0.285)};
\addplot[lineaE, flechaen=0.604] coordinates {(-0.963,0.285) (-0.999,0.302) (-1.036,0.319) (-1.072,0.337) (-1.108,0.355) (-1.144,0.374) (-1.180,0.392) (-1.216,0.411) (-1.251,0.431) (-1.287,0.450) (-1.322,0.470) (-1.357,0.490) (-1.392,0.511) (-1.427,0.531) (-1.461,0.552) (-1.496,0.573) (-1.530,0.595) (-1.564,0.617) (-1.598,0.639) (-1.632,0.661) (-1.666,0.683) (-1.699,0.706) (-1.732,0.729) (-1.766,0.752) (-1.799,0.775) (-1.831,0.799) (-1.864,0.823) (-1.897,0.847) (-1.929,0.871) (-1.961,0.895) (-1.993,0.920) (-2.025,0.945) (-2.057,0.970) (-2.089,0.995) (-2.120,1.021) (-2.151,1.047) (-2.182,1.072) (-2.213,1.098) (-2.244,1.125) (-2.274,1.151) (-2.305,1.178) (-2.335,1.205) (-2.365,1.232) (-2.395,1.259) (-2.425,1.286) (-2.455,1.314) (-2.484,1.342) (-2.513,1.370) (-2.542,1.398) (-2.571,1.426) (-2.600,1.454)};
\addplot[lineaE, flechaen=0.074] coordinates {(2.052,2.000) (2.038,1.962) (2.024,1.924) (2.009,1.887) (1.994,1.849) (1.979,1.812) (1.963,1.775) (1.946,1.738) (1.930,1.701) (1.913,1.665) (1.895,1.628) (1.877,1.592) (1.859,1.556) (1.840,1.520) (1.821,1.485) (1.801,1.450) (1.781,1.415) (1.761,1.380) (1.740,1.345) (1.719,1.311) (1.698,1.276) (1.676,1.242) (1.654,1.209) (1.631,1.175) (1.608,1.142) (1.585,1.109) (1.561,1.077) (1.537,1.044) (1.512,1.012) (1.487,0.980) (1.462,0.949) (1.437,0.918) (1.411,0.887) (1.384,0.856) (1.358,0.826) (1.330,0.796) (1.303,0.766) (1.275,0.737) (1.247,0.708) (1.219,0.680) (1.190,0.651) (1.161,0.623) (1.131,0.596) (1.101,0.569) (1.071,0.542) (1.040,0.516) (1.009,0.490) (0.978,0.464) (0.946,0.439) (0.914,0.415)};
\addplot[lineaE, flechaen=0.926] coordinates {(-0.914,0.415) (-0.946,0.439) (-0.978,0.464) (-1.009,0.490) (-1.040,0.516) (-1.071,0.542) (-1.101,0.569) (-1.131,0.596) (-1.161,0.623) (-1.190,0.651) (-1.219,0.680) (-1.247,0.708) (-1.275,0.737) (-1.303,0.766) (-1.330,0.796) (-1.358,0.826) (-1.384,0.856) (-1.411,0.887) (-1.437,0.918) (-1.462,0.949) (-1.487,0.980) (-1.512,1.012) (-1.537,1.044) (-1.561,1.077) (-1.585,1.109) (-1.608,1.142) (-1.631,1.175) (-1.654,1.209) (-1.676,1.242) (-1.698,1.276) (-1.719,1.311) (-1.740,1.345) (-1.761,1.380) (-1.781,1.415) (-1.801,1.450) (-1.821,1.485) (-1.840,1.520) (-1.859,1.556) (-1.877,1.592) (-1.895,1.628) (-1.913,1.665) (-1.930,1.701) (-1.946,1.738) (-1.963,1.775) (-1.979,1.812) (-1.994,1.849) (-2.009,1.887) (-2.024,1.924) (-2.038,1.962) (-2.052,2.000)};
\addplot[lineaE, flechaen=0.500] coordinates {(1.374,2.000) (1.375,1.959) (1.374,1.918) (1.374,1.877) (1.372,1.837) (1.369,1.796) (1.366,1.755) (1.362,1.714) (1.357,1.674) (1.351,1.633) (1.345,1.593) (1.338,1.553) (1.330,1.513) (1.321,1.473) (1.312,1.433) (1.302,1.393) (1.291,1.354) (1.279,1.315) (1.267,1.276) (1.254,1.237) (1.240,1.199) (1.226,1.160) (1.211,1.122) (1.195,1.085) (1.178,1.047) (1.161,1.010) (1.143,0.974) (1.125,0.937) (1.105,0.901) (1.085,0.865) (1.065,0.830) (1.044,0.795) (1.022,0.761) (0.999,0.727) (0.976,0.693) (0.952,0.660) (0.928,0.627) (0.903,0.595) (0.877,0.563) (0.851,0.532)};
\addplot[lineaE, flechaen=0.500] coordinates {(-0.851,0.532) (-0.877,0.563) (-0.903,0.595) (-0.928,0.627) (-0.952,0.660) (-0.976,0.693) (-0.999,0.727) (-1.022,0.761) (-1.044,0.795) (-1.065,0.830) (-1.085,0.865) (-1.105,0.901) (-1.125,0.937) (-1.143,0.974) (-1.161,1.010) (-1.178,1.047) (-1.195,1.085) (-1.211,1.122) (-1.226,1.160) (-1.240,1.199) (-1.254,1.237) (-1.267,1.276) (-1.279,1.315) (-1.291,1.354) (-1.302,1.393) (-1.312,1.433) (-1.321,1.473) (-1.330,1.513) (-1.338,1.553) (-1.345,1.593) (-1.351,1.633) (-1.357,1.674) (-1.362,1.714) (-1.366,1.755) (-1.369,1.796) (-1.372,1.837) (-1.374,1.877) (-1.374,1.918) (-1.375,1.959) (-1.374,2.000)};
\addplot[lineaE, flechaen=0.500] coordinates {(0.763,0.650) (0.780,0.687) (0.797,0.723) (0.812,0.761) (0.827,0.799) (0.840,0.837) (0.853,0.875) (0.864,0.914) (0.874,0.954) (0.883,0.993) (0.891,1.033) (0.898,1.073) (0.904,1.113) (0.909,1.153) (0.913,1.193) (0.915,1.234) (0.916,1.274) (0.916,1.315) (0.915,1.355) (0.913,1.396) (0.909,1.436) (0.904,1.476) (0.898,1.516) (0.891,1.556) (0.882,1.596) (0.873,1.635) (0.861,1.674) (0.849,1.712) (0.835,1.750) (0.820,1.788) (0.804,1.825) (0.786,1.862) (0.767,1.897) (0.747,1.932) (0.725,1.967) (0.702,2.000)};
\addplot[lineaE, flechaen=0.500] coordinates {(-0.702,2.000) (-0.725,1.967) (-0.747,1.932) (-0.767,1.897) (-0.786,1.862) (-0.804,1.825) (-0.820,1.788) (-0.835,1.750) (-0.849,1.712) (-0.861,1.674) (-0.873,1.635) (-0.882,1.596) (-0.891,1.556) (-0.898,1.516) (-0.904,1.476) (-0.909,1.436) (-0.913,1.396) (-0.915,1.355) (-0.916,1.315) (-0.916,1.274) (-0.915,1.234) (-0.913,1.193) (-0.909,1.153) (-0.904,1.113) (-0.898,1.073) (-0.891,1.033) (-0.883,0.993) (-0.874,0.954) (-0.864,0.914) (-0.853,0.875) (-0.840,0.837) (-0.827,0.799) (-0.812,0.761) (-0.797,0.723) (-0.780,0.687) (-0.763,0.650)};
\addplot[lineaE, flechaen=0.708, flechaen=0.292] coordinates {(0.633,0.775) (0.637,0.816) (0.640,0.856) (0.641,0.896) (0.641,0.937) (0.639,0.978) (0.635,1.018) (0.629,1.058) (0.622,1.098) (0.613,1.137) (0.602,1.176) (0.589,1.215) (0.574,1.253) (0.557,1.289) (0.538,1.325) (0.518,1.360) (0.496,1.394) (0.471,1.427) (0.445,1.458) (0.418,1.487) (0.388,1.515) (0.357,1.541) (0.324,1.565) (0.290,1.587) (0.254,1.606) (0.217,1.623) (0.179,1.637) (0.140,1.649) (0.101,1.657) (0.061,1.663) (0.020,1.666) (-0.020,1.666) (-0.061,1.663) (-0.101,1.657) (-0.140,1.649) (-0.179,1.637) (-0.217,1.623) (-0.254,1.606) (-0.290,1.587) (-0.324,1.565) (-0.357,1.541) (-0.388,1.515) (-0.418,1.487) (-0.445,1.458) (-0.471,1.427) (-0.496,1.394) (-0.518,1.360) (-0.538,1.325) (-0.557,1.289) (-0.574,1.253) (-0.589,1.215) (-0.602,1.176) (-0.613,1.137) (-0.622,1.098) (-0.629,1.058) (-0.635,1.018) (-0.639,0.978) (-0.641,0.937) (-0.641,0.896) (-0.640,0.856) (-0.637,0.816) (-0.633,0.775)};
\addplot[lineaE, flechaen=0.500] coordinates {(0.423,0.908) (0.406,0.946) (0.387,0.982) (0.365,1.017) (0.340,1.050) (0.313,1.081) (0.283,1.110) (0.250,1.136) (0.216,1.159) (0.180,1.178) (0.142,1.194) (0.102,1.206) (0.062,1.215) (0.021,1.219) (-0.021,1.219) (-0.062,1.215) (-0.102,1.206) (-0.142,1.194) (-0.180,1.178) (-0.216,1.159) (-0.250,1.136) (-0.283,1.110) (-0.313,1.081) (-0.340,1.050) (-0.365,1.017) (-0.387,0.982) (-0.406,0.946) (-0.423,0.908)};
\addplot[lineaE, flechaen=0.396] coordinates {(2.600,-1.454) (2.571,-1.426) (2.542,-1.398) (2.513,-1.370) (2.484,-1.342) (2.455,-1.314) (2.425,-1.286) (2.395,-1.259) (2.365,-1.232) (2.335,-1.205) (2.305,-1.178) (2.274,-1.151) (2.244,-1.125) (2.213,-1.098) (2.182,-1.072) (2.151,-1.047) (2.120,-1.021) (2.089,-0.995) (2.057,-0.970) (2.025,-0.945) (1.993,-0.920) (1.961,-0.895) (1.929,-0.871) (1.897,-0.847) (1.864,-0.823) (1.831,-0.799) (1.799,-0.775) (1.766,-0.752) (1.732,-0.729) (1.699,-0.706) (1.666,-0.683) (1.632,-0.661) (1.598,-0.639) (1.564,-0.617) (1.530,-0.595) (1.496,-0.573) (1.461,-0.552) (1.427,-0.531) (1.392,-0.511) (1.357,-0.490) (1.322,-0.470) (1.287,-0.450) (1.251,-0.431) (1.216,-0.411) (1.180,-0.392) (1.144,-0.374) (1.108,-0.355) (1.072,-0.337) (1.036,-0.319) (0.999,-0.302) (0.963,-0.285)};
\addplot[lineaE, flechaen=0.604] coordinates {(-0.963,-0.285) (-0.999,-0.302) (-1.036,-0.319) (-1.072,-0.337) (-1.108,-0.355) (-1.144,-0.374) (-1.180,-0.392) (-1.216,-0.411) (-1.251,-0.431) (-1.287,-0.450) (-1.322,-0.470) (-1.357,-0.490) (-1.392,-0.511) (-1.427,-0.531) (-1.461,-0.552) (-1.496,-0.573) (-1.530,-0.595) (-1.564,-0.617) (-1.598,-0.639) (-1.632,-0.661) (-1.666,-0.683) (-1.699,-0.706) (-1.732,-0.729) (-1.766,-0.752) (-1.799,-0.775) (-1.831,-0.799) (-1.864,-0.823) (-1.897,-0.847) (-1.929,-0.871) (-1.961,-0.895) (-1.993,-0.920) (-2.025,-0.945) (-2.057,-0.970) (-2.089,-0.995) (-2.120,-1.021) (-2.151,-1.047) (-2.182,-1.072) (-2.213,-1.098) (-2.244,-1.125) (-2.274,-1.151) (-2.305,-1.178) (-2.335,-1.205) (-2.365,-1.232) (-2.395,-1.259) (-2.425,-1.286) (-2.455,-1.314) (-2.484,-1.342) (-2.513,-1.370) (-2.542,-1.398) (-2.571,-1.426) (-2.600,-1.454)};
\addplot[lineaE, flechaen=0.074] coordinates {(2.052,-2.000) (2.038,-1.962) (2.024,-1.924) (2.009,-1.887) (1.994,-1.849) (1.979,-1.812) (1.963,-1.775) (1.946,-1.738) (1.930,-1.701) (1.913,-1.665) (1.895,-1.628) (1.877,-1.592) (1.859,-1.556) (1.840,-1.520) (1.821,-1.485) (1.801,-1.450) (1.781,-1.415) (1.761,-1.380) (1.740,-1.345) (1.719,-1.311) (1.698,-1.276) (1.676,-1.242) (1.654,-1.209) (1.631,-1.175) (1.608,-1.142) (1.585,-1.109) (1.561,-1.077) (1.537,-1.044) (1.512,-1.012) (1.487,-0.980) (1.462,-0.949) (1.437,-0.918) (1.411,-0.887) (1.384,-0.856) (1.358,-0.826) (1.330,-0.796) (1.303,-0.766) (1.275,-0.737) (1.247,-0.708) (1.219,-0.680) (1.190,-0.651) (1.161,-0.623) (1.131,-0.596) (1.101,-0.569) (1.071,-0.542) (1.040,-0.516) (1.009,-0.490) (0.978,-0.464) (0.946,-0.439) (0.914,-0.415)};
\addplot[lineaE, flechaen=0.926] coordinates {(-0.914,-0.415) (-0.946,-0.439) (-0.978,-0.464) (-1.009,-0.490) (-1.040,-0.516) (-1.071,-0.542) (-1.101,-0.569) (-1.131,-0.596) (-1.161,-0.623) (-1.190,-0.651) (-1.219,-0.680) (-1.247,-0.708) (-1.275,-0.737) (-1.303,-0.766) (-1.330,-0.796) (-1.358,-0.826) (-1.384,-0.856) (-1.411,-0.887) (-1.437,-0.918) (-1.462,-0.949) (-1.487,-0.980) (-1.512,-1.012) (-1.537,-1.044) (-1.561,-1.077) (-1.585,-1.109) (-1.608,-1.142) (-1.631,-1.175) (-1.654,-1.209) (-1.676,-1.242) (-1.698,-1.276) (-1.719,-1.311) (-1.740,-1.345) (-1.761,-1.380) (-1.781,-1.415) (-1.801,-1.450) (-1.821,-1.485) (-1.840,-1.520) (-1.859,-1.556) (-1.877,-1.592) (-1.895,-1.628) (-1.913,-1.665) (-1.930,-1.701) (-1.946,-1.738) (-1.963,-1.775) (-1.979,-1.812) (-1.994,-1.849) (-2.009,-1.887) (-2.024,-1.924) (-2.038,-1.962) (-2.052,-2.000)};
\addplot[lineaE, flechaen=0.500] coordinates {(1.374,-2.000) (1.375,-1.959) (1.374,-1.918) (1.374,-1.877) (1.372,-1.837) (1.369,-1.796) (1.366,-1.755) (1.362,-1.714) (1.357,-1.674) (1.351,-1.633) (1.345,-1.593) (1.338,-1.553) (1.330,-1.513) (1.321,-1.473) (1.312,-1.433) (1.302,-1.393) (1.291,-1.354) (1.279,-1.315) (1.267,-1.276) (1.254,-1.237) (1.240,-1.199) (1.226,-1.160) (1.211,-1.122) (1.195,-1.085) (1.178,-1.047) (1.161,-1.010) (1.143,-0.974) (1.125,-0.937) (1.105,-0.901) (1.085,-0.865) (1.065,-0.830) (1.044,-0.795) (1.022,-0.761) (0.999,-0.727) (0.976,-0.693) (0.952,-0.660) (0.928,-0.627) (0.903,-0.595) (0.877,-0.563) (0.851,-0.532)};
\addplot[lineaE, flechaen=0.500] coordinates {(-0.851,-0.532) (-0.877,-0.563) (-0.903,-0.595) (-0.928,-0.627) (-0.952,-0.660) (-0.976,-0.693) (-0.999,-0.727) (-1.022,-0.761) (-1.044,-0.795) (-1.065,-0.830) (-1.085,-0.865) (-1.105,-0.901) (-1.125,-0.937) (-1.143,-0.974) (-1.161,-1.010) (-1.178,-1.047) (-1.195,-1.085) (-1.211,-1.122) (-1.226,-1.160) (-1.240,-1.199) (-1.254,-1.237) (-1.267,-1.276) (-1.279,-1.315) (-1.291,-1.354) (-1.302,-1.393) (-1.312,-1.433) (-1.321,-1.473) (-1.330,-1.513) (-1.338,-1.553) (-1.345,-1.593) (-1.351,-1.633) (-1.357,-1.674) (-1.362,-1.714) (-1.366,-1.755) (-1.369,-1.796) (-1.372,-1.837) (-1.374,-1.877) (-1.374,-1.918) (-1.375,-1.959) (-1.374,-2.000)};
\addplot[lineaE, flechaen=0.500] coordinates {(0.763,-0.650) (0.780,-0.687) (0.797,-0.723) (0.812,-0.761) (0.827,-0.799) (0.840,-0.837) (0.853,-0.875) (0.864,-0.914) (0.874,-0.954) (0.883,-0.993) (0.891,-1.033) (0.898,-1.073) (0.904,-1.113) (0.909,-1.153) (0.913,-1.193) (0.915,-1.234) (0.916,-1.274) (0.916,-1.315) (0.915,-1.355) (0.913,-1.396) (0.909,-1.436) (0.904,-1.476) (0.898,-1.516) (0.891,-1.556) (0.882,-1.596) (0.873,-1.635) (0.861,-1.674) (0.849,-1.712) (0.835,-1.750) (0.820,-1.788) (0.804,-1.825) (0.786,-1.862) (0.767,-1.897) (0.747,-1.932) (0.725,-1.967) (0.702,-2.000)};
\addplot[lineaE, flechaen=0.500] coordinates {(-0.702,-2.000) (-0.725,-1.967) (-0.747,-1.932) (-0.767,-1.897) (-0.786,-1.862) (-0.804,-1.825) (-0.820,-1.788) (-0.835,-1.750) (-0.849,-1.712) (-0.861,-1.674) (-0.873,-1.635) (-0.882,-1.596) (-0.891,-1.556) (-0.898,-1.516) (-0.904,-1.476) (-0.909,-1.436) (-0.913,-1.396) (-0.915,-1.355) (-0.916,-1.315) (-0.916,-1.274) (-0.915,-1.234) (-0.913,-1.193) (-0.909,-1.153) (-0.904,-1.113) (-0.898,-1.073) (-0.891,-1.033) (-0.883,-0.993) (-0.874,-0.954) (-0.864,-0.914) (-0.853,-0.875) (-0.840,-0.837) (-0.827,-0.799) (-0.812,-0.761) (-0.797,-0.723) (-0.780,-0.687) (-0.763,-0.650)};
\addplot[lineaE, flechaen=0.708, flechaen=0.292] coordinates {(0.633,-0.775) (0.637,-0.816) (0.640,-0.856) (0.641,-0.896) (0.641,-0.937) (0.639,-0.978) (0.635,-1.018) (0.629,-1.058) (0.622,-1.098) (0.613,-1.137) (0.602,-1.176) (0.589,-1.215) (0.574,-1.253) (0.557,-1.289) (0.538,-1.325) (0.518,-1.360) (0.496,-1.394) (0.471,-1.427) (0.445,-1.458) (0.418,-1.487) (0.388,-1.515) (0.357,-1.541) (0.324,-1.565) (0.290,-1.587) (0.254,-1.606) (0.217,-1.623) (0.179,-1.637) (0.140,-1.649) (0.101,-1.657) (0.061,-1.663) (0.020,-1.666) (-0.020,-1.666) (-0.061,-1.663) (-0.101,-1.657) (-0.140,-1.649) (-0.179,-1.637) (-0.217,-1.623) (-0.254,-1.606) (-0.290,-1.587) (-0.324,-1.565) (-0.357,-1.541) (-0.388,-1.515) (-0.418,-1.487) (-0.445,-1.458) (-0.471,-1.427) (-0.496,-1.394) (-0.518,-1.360) (-0.538,-1.325) (-0.557,-1.289) (-0.574,-1.253) (-0.589,-1.215) (-0.602,-1.176) (-0.613,-1.137) (-0.622,-1.098) (-0.629,-1.058) (-0.635,-1.018) (-0.639,-0.978) (-0.641,-0.937) (-0.641,-0.896) (-0.640,-0.856) (-0.637,-0.816) (-0.633,-0.775)};
\addplot[lineaE, flechaen=0.500] coordinates {(0.423,-0.908) (0.406,-0.946) (0.387,-0.982) (0.365,-1.017) (0.340,-1.050) (0.313,-1.081) (0.283,-1.110) (0.250,-1.136) (0.216,-1.159) (0.180,-1.178) (0.142,-1.194) (0.102,-1.206) (0.062,-1.215) (0.021,-1.219) (-0.021,-1.219) (-0.062,-1.215) (-0.102,-1.206) (-0.142,-1.194) (-0.180,-1.178) (-0.216,-1.159) (-0.250,-1.136) (-0.283,-1.110) (-0.313,-1.081) (-0.340,-1.050) (-0.365,-1.017) (-0.387,-0.982) (-0.406,-0.946) (-0.423,-0.908)};
\addplot[lineaE, flechaen=0.625] coordinates {(1.000,0.000) (2.600,0.000)};
\addplot[lineaE, flechaen=0.375] coordinates {(-2.600,0.000) (-1.000,0.000)};

  \draw[texto, line width=0.6pt, dashed] (axis cs:0,0) circle[radius=1];
  \draw[vecP] (axis cs:-0.32,0) -- (axis cs:0.32,0);
  \node[etiqm, font=\footnotesize] at (axis cs:0,0.3) {$\vp$};
  \node[etiqm, font=\footnotesize] at (axis cs:-0.45,-0.55) {$r<a$};
\end{axis}
\end{tikzpicture}
\caption{Campo del dipolo para $r>a$}
\end{subfigure}\hfill
\begin{subfigure}[b]{0.44\textwidth}\centering
\begin{tikzpicture}
\begin{polaraxis}[width=5.6cm, xtick={0,45,90,135,180,225,270,315}, ytick={2,4}, yticklabels={}, ymax=4.3,
   xticklabel={$\pgfmathprintnumber{\tick}^\circ$}, tick label style={font=\tiny},
   grid style={draw=gris, line width=0.3pt}, axis line style={draw=gris}]
  \addplot[curvaE, domain=0:360, samples=181, fill=campoE, fill opacity=0.15] {3*cos(x)^2+1};
\end{polaraxis}
\end{tikzpicture}
\caption{$u_e(r,\theta)\propto 3\cos^2\theta+1$}
\end{subfigure}
\caption{\Cref{prob:dipolo}. (a) Líneas del campo de un dipolo puntual $\vp$ (horizontal) fuera de la esfera de radio $a$; el interior,
donde el campo diverge al acercarse al dipolo, queda excluido del cálculo. (b) A distancia fija, la densidad de energía es cuatro
veces mayor en el eje del dipolo ($\theta=0,\pi$) que en el plano ecuatorial ($\theta=\pi/2$), donde el campo es la mitad.}
\label{fig:p52}
\end{figure}
